\section{Results}

\subsection{Competitive District Fraction}

\begin{table}[h]
\centering\small
\caption{Competitive district fraction and efficiency gap.
Bisect vs.\ enacted 2022 plans.
Competitive = margin $< 10$pp; EG computed from VEST 2020 presidential.}
\label{tab:competitiveness}
\begin{tabular}{lcccccc}
\toprule
& \multicolumn{2}{c}{\textbf{Bisect}} & \multicolumn{2}{c}{\textbf{Enacted}} & \multicolumn{2}{c}{\textbf{Difference}} \\
State & Comp.\% & EG & Comp.\% & EG & $\Delta$Comp. & $\Delta$EG \\
\midrule
NC  ($k=14$) & 50\% & $-3.1\%$ & 14\% & $+6.2\%$ & $+36$pp & $-9.3$pp \\
WI  ($k=8$)  & 38\% & $-2.5\%$ & 13\% & $+7.8\%$ & $+25$pp & $-10.3$pp \\
TX  ($k=38$) & 34\% & $-2.8\%$ & 16\% & $+4.3\%$ & $+18$pp & $-7.1$pp \\
FL  ($k=28$) & 39\% & $-3.2\%$ & 18\% & $+5.4\%$ & $+21$pp & $-8.6$pp \\
PA  ($k=17$) & 41\% & $-2.9\%$ & 24\% & $+4.8\%$ & $+17$pp & $-7.7$pp \\
\midrule
\textbf{Mean} & \textbf{40\%} & $\mathbf{-2.9\%}$ & \textbf{17\%} & $\mathbf{+5.7\%}$ & $\mathbf{+23}$\textbf{pp} & $\mathbf{-8.6}$\textbf{pp} \\
\bottomrule
\end{tabular}
\end{table}

\textsc{Bisect} plans produce 40\% competitive districts on average vs.\ 17\%
for enacted gerrymanders --- a 23 percentage point increase.
The competitive district fraction under \textsc{bisect} approaches the fraction
in Iowa (44\% competitive, 2022 cycle), the established benchmark for
neutral redistricting.

The efficiency gap for \textsc{bisect} plans ($-2.9\%$) is close to zero
and reflects the geographic sorting of Democratic voters into urban cores.
The enacted plans' EG ($+5.7\%$) significantly amplifies the baseline
geographic advantage for Republicans beyond what geography alone produces.

\subsection{Geographic Explanation}

The competitive fraction difference arises from two sources:
(1) \textsc{bisect} cannot pack opposition voters, so districts near the
urban-suburban boundary --- which tend to be competitive --- remain whole
rather than being cracked across multiple districts;
(2) compactness optimisation draws districts along geographic gradients
(urban core, inner suburb, outer suburb, exurban) that correspond to
partisan gradients, creating ``staircase'' districts with intermediate
partisan compositions.
